\section{Background}\label{sec:background}

\subsection{Chips, Slices, and Hosts}\label{sec:bg-tpu}
A TPU v5e chip has 16~GiB of high-bandwidth memory (HBM)~\cite{googletpuv5e}. % src: googletpuv5e as read on 2026-11-23 ("16 GiB"), paper-v4/SOURCE_NOTES.md section 1
Google's TPUs link their chips with an inter-chip interconnect (ICI)~\cite{jouppi2020domain}.
The full TPU v5e configuration of 256 chips has a two-dimensional torus topology~\cite{googletpuv5e}. % src: googletpuv5e as read on 2026-11-23 ("TPU Pod size 256 chips", "2D torus")
Our slice had eight chips in a 2~$\times$~4 topology. % src: configs/af_spike_job.yaml -> nodeSelector; recomputed: V2.topology
A host is the CPU machine that runs the Python process and drives the chips. JAX treats each chip as one device of that
process. In this paper, a chip is one TPU v5e chip. A slice is the set of eight chips that each Kubernetes Job received
(Section~\ref{sec:platforms}).

\subsection{Compilation and Mapping in JAX}\label{sec:bg-jax}
Before compiling a function, JAX traces the function into an intermediate program. \texttt{jax.jit} compiles a function
with XLA when the function is first called, caches the compiled code, and reuses that code in later
calls~\cite{jaxjit}. In this paper, compilation means tracing and XLA compilation together. We do not separate the two.
JAX can also store compiled programs on disk, so that a later run can load them instead of compiling
again~\cite{jaxcache}. We call these stored programs the compilation cache.

JAX offers several ways to run more than one query. \texttt{jax.vmap} vectorizes a function over a batch axis of its
inputs~\cite{jaxvmap}. \texttt{jax.vmap} does not by itself spread the batch over devices. \texttt{jax.pmap} replicates
a function and runs each replica on its own device. Each device receives one element of the leading axis of the
input~\cite{jaxpmap}. The documentation of \texttt{jax.pmap} now recommends \texttt{jax.shard\_map}
instead~\cite{jaxpmap}.
Since JAX 0.10.0, \texttt{jax.pmap} is implemented as \texttt{jax.jit} applied to \texttt{jax.shard\_map}~\cite{jaxchangelog}. % src: jaxchangelog ("JAX 0.10.0", "jax.pmap will now always use the new implementation that wraps jax.jit(jax.shard_map)")
Our external ensemble (Section~\ref{sec:res-auto}) uses \texttt{jax.lax.pmean}, which computes the mean of an array over
the devices of a mapped axis~\cite{jaxpmean}.

JAX can also partition one computation automatically over a mesh. A mesh is a grid of devices with named axes. The
GSPMD partitioner of XLA starts from the sharding annotations that the program places on some arrays. The partitioner
derives the placement of the other arrays from these annotations~\cite{xu2021gspmd}. The OpenXLA guide to Shardy for
JAX users says that JAX now always uses the Shardy partitioner~\cite{shardy2026}.
We read this guide on November 23, 2026. % src: shardy2026 ("read 2026-11-23", "Shardy is now always enabled in JAX"), paper-v4/SOURCE_NOTES.md section 2
According to the JAX change log, Shardy became the default partitioner in JAX 0.7.0, and JAX 0.12.1 removed the GSPMD partitioner~\cite{jaxchangelog}. % src: jaxchangelog ("JAX 0.7.0", "JAX is migrating from GSPMD to Shardy by default", "JAX 0.12.1", "Removed support for the legacy GSPMD partitioner")
The committed Jobs installed JAX 0.10.2 and set no partitioner flag, so Shardy was the default. % src: recomputed: V2.jax_pin; src/spike_meshshard_forward_pass.py ("Job pins jax[tpu]==0.10.2, which defaults to Shardy")
The runs did not record which partitioner ran. In Shardy, an array is replicated along each mesh axis that does not
shard one of the dimensions of the array~\cite{shardyrepr}. This rule matters for our partitioning experiment, which
placed no sharding annotations (Section~\ref{sec:res-auto}).

\subsection{AlphaFold2 Inference Pipeline}\label{sec:bg-af2}
AlphaFold2 reads a multiple sequence alignment (MSA). An MSA lines up the query with related protein sequences. A
stack of Evoformer blocks updates a representation of the MSA and a representation of residue pairs. Residues are the
amino acids of the protein chain. The model also reads extra MSA rows, a second and larger set of aligned sequences.
Before the 48 Evoformer blocks, a separate stack of four smaller blocks processes these extra rows~\cite{alphafoldrepo}. % src: AlphaFold2 at commit c77e5d2, alphafold/model/config.py:240 ("'evoformer_num_block': 48"), config.py:307-308 ("'extra_msa_channel': 64", "'extra_msa_stack_num_block': 4"), config.py:"msa_channel": 256
A structure module then turns these representations into atom coordinates~\cite{jumper2021alphafold}. The public
implementation uses Haiku~\cite{haiku2020github}. Haiku is a neural-network library for JAX. The model can recycle its
output as a new input for a fixed number of iterations. The public implementation runs recycling in a
\texttt{while\_loop}. By default, the monomer presets of the public implementation run three recycling iterations and
predict each target with five models, \texttt{model\_1} to \texttt{model\_5}~\cite{alphafoldrepo}. % src: AlphaFold2 at commit c77e5d2, alphafold/model/config.py:146 ("'num_recycle': 3"), config.py:30-37 (MODEL_PRESETS 'monomer': model_1 to model_5)

The implementation can also run an ensemble. The members of an ensemble are differently featurized copies of the same
input. Featurization turns the sequence and the MSA into the arrays that the model reads. The implementation computes the
first member outside a loop and the others in a sequential Haiku \texttt{while\_loop}~\cite{alphafoldrepo}. % src: AlphaFold2 at commit c77e5d2, alphafold/model/modules.py:174-205 ("hk.while_loop"), paper-v4/SOURCE_NOTES.md section 3; recomputed: SEP.cpu.0915.alphafold_commit
The authors of AlphaFold2 turn off ensembling for most inference~\cite{jumper2021alphafold}. In the released code, the
preset \texttt{monomer\_casp14} uses eight ensemble members, and the other presets use one~\cite{alphafoldrepo}. % src: AlphaFold2 at commit c77e5d2, run_alphafold.py:601-604 ("num_ensemble = 8", "num_ensemble = 1"); jumper2021alphafold ("we turn off ensembling for most inference")
Our runs used one ensemble member, so the ensemble \texttt{while\_loop} ran no iteration in the timed calls. % src: src/spike_tpu_forward_pass.py ("cfg.data.eval.num_ensemble = 1")

Our scripts time one of two methods of the model. The \texttt{apply} method is the forward pass on the device. The
\texttt{predict} method runs \texttt{apply} and then computes confidence metrics on the host with NumPy.
